\chapter{Guía de Referencia Rápida de \fenics}
\label{app:fenics}

Este apéndice proporciona una referencia rápida de los comandos y conceptos más comunes de \fenics utilizados en este libro. No pretende ser exhaustivo, sino servir como un recordatorio práctico para la sintaxis clave.

%==================================================================================
\section{Flujo de Trabajo Típico de un Script}
%==================================================================================
Un script de \fenics para un problema lineal sigue un patrón de 5 pasos bien definidos.
\begin{enumerate}
    \item \textbf{Malla y Espacio Funcional:} Definir el dominio computacional y el tipo de aproximación para la solución.
    \item \textbf{Condiciones de Contorno:} Especificar los valores conocidos de la solución en la frontera (Dirichlet).
    \item \textbf{Problema Variacional:} Traducir la forma débil \(a(u,v) = L(v)\) a código.
    \item \textbf{Resolver:} Invocar al solver para encontrar la función de solución.
    \item \textbf{Post-procesamiento:} Guardar o visualizar los resultados.
\end{enumerate}

%==================================================================================
\section{Comandos Esenciales}
%==================================================================================

%----------------------------------------------------------------------------------
\subsection{Malla y Geometría}
%----------------------------------------------------------------------------------
\begin{minted}[frame=lines, fontsize=\small]{python}
from fenics import *
from mshr import * # Para geometrías más complejas

# Malla de un cuadrado unitario de 32x32 elementos
mesh = UnitSquareMesh(32, 32)

# Malla de un rectángulo
mesh = RectangleMesh(Point(0,0), Point(L, H), nx, ny)

# Malla de un disco
domain = Circle(Point(cx, cy), radius)
mesh = generate_mesh(domain, resolution)
\end{minted}

%----------------------------------------------------------------------------------
\subsection{Espacio de Funciones}
%----------------------------------------------------------------------------------
\begin{minted}[frame=lines, fontsize=\small]{python}
# Espacio de funciones escalares (para Temperatura, Presión)
# Elementos Lagrange de primer orden (lineales)
P1 = FiniteElement("P", mesh.ufl_cell(), 1)
# Elementos Lagrange de segundo orden (cuadráticos)
P2 = FiniteElement("P", mesh.ufl_cell(), 2)

# Espacio de funciones vectoriales (para Desplazamiento, Velocidad)
V_P2 = VectorElement("P", mesh.ufl_cell(), 2)

# Creación del FunctionSpace
V = FunctionSpace(mesh, P1)
W = VectorFunctionSpace(mesh, "P", 2) # Atajo para vectoriales

# Espacio mixto para Taylor-Hood (Velocidad-Presión)
ME = MixedElement([V_P2, P1])
W_mixed = FunctionSpace(mesh, ME)
\end{minted}

%----------------------------------------------------------------------------------
\subsection{Definición de Funciones}
%----------------------------------------------------------------------------------
\begin{minted}[frame=lines, fontsize=\small]{python}
# Funciones para la forma débil
u = TrialFunction(V) # Incógnita principal, ej. T o u
v = TestFunction(V)  # Función de prueba

# Función para almacenar una solución calculada
u_sol = Function(V)

# Función del paso de tiempo anterior
u_n = Function(V)

# Constantes y Expresiones
c = Constant(5.0)
f = Expression("sin(x[0])*cos(x[1])", degree=2)
\end{minted}

%----------------------------------------------------------------------------------
\subsection{Condiciones de Contorno (Dirichlet)}
%----------------------------------------------------------------------------------
\begin{minted}[frame=lines, fontsize=\small]{python}
# 1. Definir el valor en el contorno
u_D = Constant(0.0) # Para un valor constante
u_expr = Expression("1 + x[0]*x[0]", degree=2) # Para un valor variable

# 2. Definir la localización del contorno (toda la frontera)
def boundary(x, on_boundary):
    return on_boundary

# 3. Crear el objeto DirichletBC
bc = DirichletBC(V, u_D, boundary)
\end{minted}

%----------------------------------------------------------------------------------
\subsection{Problema Variacional y Solución}
%----------------------------------------------------------------------------------
\paragraph{Problemas Lineales}
\begin{minted}[frame=lines, fontsize=\small]{python}
# Forma débil: a(u,v) = L(v)
a = dot(grad(u), grad(v))*dx
L = f*v*dx

# Solución
u_sol = Function(V)
solve(a == L, u_sol, bc)
\end{minted}

\paragraph{Problemas No Lineales}
\begin{minted}[frame=lines, fontsize=\small]{python}
# La incógnita es una Function, no TrialFunction
u = Function(V)

# Definir la forma débil como un residuo F(u; v) = 0
k = 1 + u**2 # Ejemplo de no linealidad
F = k*dot(grad(u), grad(v))*dx - f*v*dx

# Solución (FEniCS invoca a Newton automáticamente)
solve(F == 0, u, bc)
\end{minted}

\paragraph{Problemas Transitorios (Bucle Temporal)}
\begin{minted}[frame=lines, fontsize=\small]{python}
for n in range(num_steps):
    # Resolver el problema para el paso actual
    solve(a == L, u_sol, bc)
    
    # Actualizar la solución del paso anterior
    u_n.assign(u_sol)
\end{-minted}

%----------------------------------------------------------------------------------
\subsection{Post-Procesamiento}
%----------------------------------------------------------------------------------
\begin{minted}[frame=lines, fontsize=\small]{python}
# Guardar en formato VTK para ParaView
vtkfile = File("resultado.pvd")
vtkfile << u_sol

# Graficar con Matplotlib (requiere 'import matplotlib.pyplot as plt')
plot(u_sol)
plt.show()

# Proyectar un tensor a un campo escalar para visualización
V_scalar = FunctionSpace(mesh, "P", 1)
sigma_xx = project(sigma(u_sol)[0,0], V_scalar)
plot(sigma_xx)
\end{minted}